% §18.3 Confidentiality Measures
%
% Monograph register: the server's leakage model, PPV as confidentiality,
% NPV/F1, multi-query leakage (prop:multi-query-leakage), worked example.
% REVISED 2026-10: the old proposition (posterior factorizes as a Kronecker
% product; I_m = m I_1; C_m = Theta(m^-1/2)) was unsupported and is replaced
% by subadditivity, the frozen-structure case, and the geometric decay of
% cover across independently keyed instances, with proof.
% Tightened; all numbers verified. Cross-refs sec:kronecker, thm:kronecker,
% def:ppv, def:npv, def:f1, sec:derived-measures, sec:trapdoor-paradigm,
% sec:maximizing-confidentiality, def:per-document-bloom, sec:union-of-sets.
% DEFINED: sec:confidentiality-measures, prop:multi-query-leakage.

\section{Confidentiality Measures}
\label{sec:confidentiality-measures}

The cipher-map framing of \Cref{sec:bloom-based} fixed an error profile
$(\fprate, 0)$ for both constructions, and
\Cref{sec:maximizing-confidentiality} worked out confidentiality at fixed
rates and prevalence. This section reads that derivation in the
encrypted-search setting: under the observation model stated there, each
chapter 7 measure has a leakage interpretation, and trading rate budget
for cover reduces to choosing $\fprate$ on the chapter 7 curve. The other
leakage channels of the setting (query equality, access pattern, result
volume) are outside this accounting.

\subsection*{The server's leakage model}

A deployed system gives the server two streams. The \emph{structural}
stream is the encrypted filter store, opaque by construction: each
$\mathsf{Enc}(B_i)$ is semantically secure, leaking nothing beyond the
filter size. The \emph{operational} stream is where the analysis lives:
each match decision is a yes/no answer to ``does $\tau_w$ correspond to a
keyword in $d_i$?'', which the server computes and therefore knows. The
plaintext keyword and document content stay hidden, but the bare fact of
the match is unavoidably revealed. Under the Bernoulli model each match
decision is a noisy bit, $1$ when the keyword truly belongs (the filter's
$\fnrate = 0$) and $1$ with probability $\fprate$ when it does not, so the
chapter 7 measures over noisy bits are the right measures for the server's
posterior.

\subsection*{PPV is the confidentiality measure}

Fix a query token $\tau_w$ and a document $d_i$ whose filter matched. The
server's posterior that $w$ truly appears in $d_i$ is, by Bayes' rule,
the positive predictive value of \Cref{def:ppv} at the cipher map's rates
and the workload prevalence $\pi$,
\[
  \Pr[w \in d_i \mid \tau_w \text{ matched}]
  = \PPV = \frac{\pi(1 - \fnrate)}{\pi(1 - \fnrate) + (1 - \pi)\fprate}.
\]
The chapter 7 formula and the server's posterior are the same
quantity derived twice.\index{positive predictive value} Following
\Cref{sec:maximizing-confidentiality}, confidentiality is the complement
$C = 1 - \PPV$, high when most matches are spurious and low when most are
genuine. At $\fnrate = 0$, which both constructions satisfy,
\[
  C \;=\; \frac{(1 - \pi)\fprate}{\pi + (1 - \pi)\fprate},
\]
the most useful equation in the chapter: on the perfect-filter
line the cover is monotone in the rate-prevalence product $(1 -
\pi)\fprate$ relative to $\pi$. The companion measure for $\bot$-reports
is the negative predictive value\index{negative predictive value} (\Cref{def:npv}), which at $\fnrate = 0$
collapses to $\NPV = 1$: a Bloom-based search reports $\bot$ only when the
keyword is truly absent (the union analysis of \Cref{sec:union-of-sets}),
so a $\bot$-report is fully reliable and the $\fnrate = 0$ regime
concentrates all confidentiality on the $\top$-side. The balanced-workload
summary is the $\Fone$-score\index{F1 score@$F_1$ score} of \Cref{def:f1}, which at $\fnrate = 0$
(recall $= 1$, precision $= \PPV$) reduces to $2\PPV/(1 + \PPV)$, monotone
in PPV and the right measure when $\pi$ is not known in advance. The chapter 7 apparatus (\Cref{sec:derived-measures}) thus carries over
unchanged, every derived measure acquiring a reading under the stated
observation model.

\subsection*{Multi-query leakage}

A single query leaks a bounded amount about one $(w, d_i)$ pair; a long
sequence leaks about the whole corpus. What accumulates, and how fast,
depends on whether a new observation is a fresh draw.

\begin{proposition}[Leakage across queries and instances]
\label{prop:multi-query-leakage}
Consider a system with profile $(\fprate, 0)$ at prevalence $\pi$.
\begin{enumerate}
\item \emph{Subadditivity.} If observations $Y_1, \ldots, Y_m$ are
  conditionally independent given the plaintext corpus $D$, then
  $I(D; Y_1, \ldots, Y_m) \le \sum_{j} I(D; Y_j)$, so per-query leakage
  bounds add.
\item \emph{Frozen structure.} Repeating a query for the pair $(w, d_i)$
  against the same structure returns the same result, so it adds no
  information about that pair (it does reveal that the query repeated).
\item \emph{Independent instances.} If the pair is tested against $t$
  structures built with independent keys over the same corpus, a true
  pair matches in all of them and a false pair in each independently with
  probability $\fprate$, so the posterior after $t$ matches is
  \[
    \PPV_t \;=\; \frac{\pi}{\pi + (1 - \pi)\fprate^{\,t}},
    \qquad
    C_t \;=\; 1 - \PPV_t \;\le\; \frac{1 - \pi}{\pi}\,\fprate^{\,t}.
  \]
\end{enumerate}
\end{proposition}

\begin{proof}
(1) $I(D; Y^m) = H(Y^m) - H(Y^m \mid D) = H(Y^m) - \sum_j H(Y_j \mid D)
\le \sum_j H(Y_j) - \sum_j H(Y_j \mid D)$, using conditional independence
for the middle equality and subadditivity of entropy for the inequality.
(2) The structure and the query token are fixed, so the result is a
deterministic function of them. (3) The $t$ match events of a false pair
are independent with probability $\fprate$ each, those of a true pair all
occur since $\fnrate = 0$; Bayes' rule gives $\PPV_t$, and dropping the
$(1 - \pi)\fprate^{\,t}$ term from the denominator of $C_t$ gives the
bound.
\end{proof}

The three parts separate what accumulates. Across distinct pairs,
leakage adds, so a fixed per-query bound gives a linear budget. Against
one structure, repetition reveals the search pattern but not more about
the pair. Across independently keyed instances, the cover of a single
fact decays geometrically: rebuilding the index under a fresh key resets
the linkage between old and new tokens, but if the observer keeps the old
index, every rebuild is another instance in part 3. Key rotation
protects only against an observer who cannot retain old structures.
Part 3 is the membership form of the coincidence attack on retrieval
codecs. There a non-member decodes, in each of $t$ independently keyed
instances, to a value drawn from the codec's distribution $(q_v)$, while
a member decodes to its true value every time; an attacker who calls an
item a member when all $t$ decodings agree is right with probability
$1 - \tfrac12 \sum_v q_v^t$ at a balanced prior. Large $\sum_v q_v^t$
means a weak attacker, so codecs concentrated on a heavy value (large
$\max_v q_v$) defend best and the uniform dense codec defends worst
(\textcite{towell-cipher-maps}; companion experiment
\texttt{bernoulli-py/\allowbreak experiments/\allowbreak 14\_coincidence\_oracle.py}).\index{codec}\index{coincidence attack}

\subsection*{A worked example}

Take the per-document construction (\Cref{def:per-document-bloom}) at
$\fprate = 0.01$, $\pi = 0.1$ (one queried keyword in ten truly appears in
the queried document). The PPV formula at $\fnrate = 0$ gives
\[
  \PPV = \frac{0.1}{0.1 + 0.9 \cdot 0.01} = \frac{0.1}{0.109}
  \approx 0.917, \qquad C \approx 0.083,
\]
a correctness-tuned, privacy-leaky point: a match gives the adversary a
$91.7\%$ posterior. Loosening to $\fprate = 0.3$ gives $\PPV = 0.1/0.37
\approx 0.270$, $C \approx 0.730$, the matches now mostly false positives
and the plaintext well protected, at the cost of a much noisier response
the client must post-filter. Every $\fprate$ at fixed $\pi$ sits on the
curve $C = (1 - \pi)\fprate / [\pi + (1 - \pi)\fprate]$; intermediate
points are $\fprate = 0.05$ at $C \approx 0.310$, $\fprate = 0.1$ at $C
\approx 0.474$, $\fprate = 0.2$ at $C \approx 0.643$. Inverting for the
design target, $\fprate^* = \pi C / [(1 - \pi)(1 - C)]$, so $C = 0.9$ at
$\pi = 0.1$ needs $\fprate^* = 1$ (degenerate), while the same target at
$\pi = 0.01$ needs $\fprate^* \approx 0.091$, a friendly operating point:
lower-prevalence workloads buy confidentiality more cheaply, matching the
chapter 17 intuition that confidentiality is abundant where the prior
already concentrates the plaintext on $\bot$. The next section turns to
the harder problem of estimating $\pi$ when the workload is adversarial.
